\section{诉讼程序}
\begin{enumerate}
  \item 有 N 多内容同民诉，同民诉的不考，考不同 \yjt 受案范围、管辖、审理期限、复议执行。
  \item 起诉。
        \begin{enumerate}
          \item 条件。同民诉。原告适格、有明确的被告、具体的诉讼请求、属于法院主管 + 管辖。
          \item 重复起诉。一事不再理。做题注意出现了多少个行为。
          \item 起诉期。
                \begin{enumerate}
                  \item 直接起诉。\xydy 6 个月。
                  \item 经过复议再起诉 15 日内 \hz 不理不睬等 2 月。
                        对不履行法定职责提起诉讼的，应当在行政机关履行法定职责期限届满之日起 6 个月内提出。
                        \begin{kk}{收到}
                          \begin{enumerate}
                            \item 全知道。知内容 + 知诉权。正常处理 \yjt 6 个月。
                            \item 知一半。内容 1 年，诉 6 月。知内容 1 年，知诉权 6 个月。同时起算，取交集。
                            \item 全不知。最长保护 5、20。最长保护期。普通 5 年，不动产 20 年。作出之日起算。
                          \end{enumerate}
                          \tcblower
                          注意。得是当事人知。
                          \al{甲乙开办丙公司。丙以土地抵押给丁，后丁办理了权属变更。
                            08 年 3 月，\xhx{甲查询}土地抵押登记情况，\xhx{发现}已变更到丁名下 \yjt 于丙而言，全不知}。
                        \end{kk}
                  \item 过起诉期不予立案。
                \end{enumerate}
          \item 仅被限制人身自由的可以口头起诉、委托。
        \end{enumerate}
  \item 受理（立案）。
        \begin{enumerate}
          \item 登记立案、裁定不予立案、先予立案（7 日内仍做不出决定）。
          \item 缺材料一次性告知 \yjt 符合/不符合 \yjt 立案/退回 \dy 裁定不予立案 \yjt 上诉，胜诉 \yjt 指令立案。
          \item 裁定驳回起诉 \yjt 上诉，胜诉 \yjt 裁定继续审理。
        \end{enumerate}
  \item 审理。
        \begin{enumerate}
          \item 基本原则。
                \begin{enumerate}
                  \item 不诉不理。
                        \begin{enumerate}
                          \item 非法院主管。\al{请求赔偿并要求\xhx{对工作人员进行处分}}。
                          \item 诉哪个审哪个判哪个。\al{先强制措施，后又进行了处罚，只\xhx{起诉处罚}}。
                        \end{enumerate}
                  \item 两审终审。除被告都能上诉（对被告是一审终审制）。
                \end{enumerate}
          \item 庭前程序，提前 3 日告知。
                \begin{enumerate}
                  \item 被告不提交答辩状不影响审理。15 日，同时也是举证期。
                  \item 增加请求。
                        \begin{enumerate}
                          \item 被告收到起诉状副本后，不得增加请求。例外。赔偿请求（最迟一审判决前，法院自主决定）。
                          \item 民事请求、其他请求。原则开庭前提出，最迟法庭调查中 + 有正当理由（同附带性审查）。
                        \end{enumerate}
                  \item 传唤当事人用传票。其他人叫通知书。
                \end{enumerate}
          \item 审理对象。
                \begin{enumerate}
                  \item 被诉行政行为。\al{公安局对甲处以 10 日拘留，甲\xhx{不服处罚}起诉。\al{市政府征收土地的决定不属于本案的审查范围 \dui}}。
                  \item 复议维持 \yjt 原行为 + 复议维持。\al{法院在对县规划局的行为作出裁判时，应对市规划局的复议决定作出裁判 \dui}。
                \end{enumerate}
          \item 审理方式。
                \begin{enumerate}
                  \item 必须开庭审理、一般能够旁听、最终判决要公开。
                  \item 政府信息公开案件是否公开法院裁断。
                \end{enumerate}
          \item 原告拒绝陈述 \yjt 释明 \yjt 承担不利后果。
          \item 审期。立案到宣判 6 个月。鉴定、勘验、公告、管辖异；调解期间、材在途不包括。常识判断。
        \end{enumerate}

        \begin{kk}{简易程序}
          \begin{enumerate}
            \item 分为法定简易、约定简易。选项说可以用简易程序请保持警惕！！
            \item 适用条件。背。
                  \begin{enumerate}
                    \item 事、争、权、义为前提；一审、级别无要求。事实清晰、争议不大、权利义务明确。
                    \item 信息公开、当场做；双方约定、钱 2 千。
                    \item 通知方式可从简，裁判文书要严肃。
                    \item 举证期限可协商，不超 15 日法院准。
                    \item 裁定转化知双方，审期仍从立案算。刑诉是转化之日起算。
                    \item 独任 45 不得延，择期判。
                  \end{enumerate}
            \item 只有处罚、诉讼有简易程序。\al{作出\xhx{冻结}的时候，可以适用简易程序 \cuo}。
            \item \al{当事人各方同意适用简易程序的，\xhx{可适用}简易程序 \dui}。
            \item \al{法官对本案\xhx{不能}适用简易程序审判 \cuo}。
            \item \al{对于甲对第一次申请信息公开答复不服提起的诉讼，法院\xhx{不能适用}简易程序审理 \cuo}。
          \end{enumerate}
        \end{kk}
  \item 二审。
        \begin{enumerate}
          \item 被告是可上诉的。
          \item 上诉状给谁都行，给二审，二审 5 日内转给一审法院。
          \item 一审判决和被诉行政行为全面审查。民诉 \yjt 当事人的上诉范围。刑诉 \yjt 不限于上诉范围。
                \al{二审认为需要改变一审判决的，应同时对县政府的决定作出判决 \dui}。
          \item 合议庭、原则开庭、3 个月。\al{二审法院\xhx{不得以}不开庭方式审理该上诉案件 \cuo}。
          \item 判决。考遗漏请求。
                \begin{enumerate}
                  \item 遗漏赔偿请求，二审认为不应赔/要赔 \yjt 欺软/怕硬 \yjt 驳回赔偿请求/赔偿部分，调解不成发回。
                  \item 遗漏普通请求。不同民诉！！
                        \begin{enumerate}
                          \item 应当裁定撤销一审判决。
                          \item 应当裁定发回一审法院重审，一审法院应当重新组成合议庭。
                        \end{enumerate}
                \end{enumerate}
        \end{enumerate}
  \item 特别制度。比较爱考。
        \begin{enumerate}
          \item 撤诉。不来、先走、不交钱 \yjt 视为撤诉。
                \begin{kk}{一改二撤三裁}
                  \begin{enumerate}
                    \item 一改。改变原行政行为，不影响继续审理。
                          \begin{enumerate}
                            \item 有错必纠。一审、二审、再审都可以改变。
                            \item 实质改变。事实、证据、行为、结果、法律适用。
                                  \al{更正处罚决定中不影响决定性质和内容的文字错误 \cuo}。
                            \item 视为改变。
                                  \begin{itemize}
                                    \item 根据原告的请求依法履行法定职责。
                                    \item 采取相应的补救、赔偿等措施。
                                    \item 在行政裁决案件中，书面认可原告与第三人达成的和解。
                                  \end{itemize}
                          \end{enumerate}
                    \item 二撤。撤回起诉。
                    \item 三裁。意思表示真实、未违反法律规定、公序良俗、第三人利益 \yjt 裁定准予撤诉。
                  \end{enumerate}
                  \tcblower
                  \begin{enumerate}
                    \item 新的行为审没审看诉没诉。另行起诉、合并审理。
                    \item 没有后悔药、撤诉裁定错了走再审。
                  \end{enumerate}
                \end{kk}
          \item 针对被告的特别制度。
                \begin{enumerate}
                  \item 缺席判决，提交的证据不能作为定案依据。
                  \item 负责人原则上要出庭。
                        \begin{enumerate}
                          \item \dydy 1 名工作人员（\al{律师 \cuo}）。
                          \item 不来不能传唤！！（不是诉讼当事人） \yjt 裁判文书写明情况、提司法建议。
                                \al{传唤其出庭 \cuo}。
                        \end{enumerate}
                \end{enumerate}
        \end{enumerate}
  \item 调解。
        \begin{enumerate}
          \item 仅针对赔偿、补偿等裁量行为，羁束不能调。处罚裁量，许可羁束。
                \begin{enumerate}
                  \item \al{罚款；税率、税额、资格证}。
                  \item \al{市监局确定的罚款数额需要水果店店主协商同意 \cuo{协商可以，需要你同意咋可能}}。
                \end{enumerate}
          \item 一般不公开、法院盖章 + 当事人签字生效、可强制执行。
          \item 调解不成应及时判决。调解协议不能成为判决的依据。有调没判，有判没调。
        \end{enumerate}
  \item 先予执行。可怜人申请可怜钱，可以复议一次。\al{行政复议 \cuo{诉讼中的复议都是司法复议}}。
  \item 庭审障碍。训诫、责令具结悔过、罚款拘留（院长批） \yjt 不服复议。选项只写复议，默认 = 司法复议。
  \item 与其他诉讼交叉。
        \begin{enumerate}
          \item 刑行。有/没有关联性 \yjt 移送且中止（等结果）/但不中止。
          \item 民行。应当分别立案，可以合并审理，应当分别裁判。\yi 分别审理。谁先谁后具体分析。\er 合并。
          \item 行政裁决附带民事。合并立、合并审、合并判。
                \al{甲乙两村就土地权属发生争议，请县政府解决。
                  县政府经查后认为土地应属于甲村，乙村不服提起行政诉讼，请求撤销县政府决定并确定土地权属}。
          \item 合议庭结论。应当另行组成合议庭，不管哪种情况都错。
          \item 是否合并，经当事人同意。\al{法院\xhx{可以决定}民事与行政案件合并审理 \cuo}。
          \item 分别收费。
        \end{enumerate}
\end{enumerate}
